\documentclass[11pt]{article}
\usepackage[T1]{fontenc}
\usepackage[utf8]{inputenc}
\usepackage{lmodern,amsmath,amssymb,amsthm,mathtools,mathrsfs,booktabs}
\usepackage[margin=1in]{geometry}
\usepackage{xurl}
\usepackage[colorlinks=true,linkcolor=blue,citecolor=blue,urlcolor=blue]{hyperref}
\hypersetup{pdftitle={Rotational linearization obstructions in a doubled Einstein--scalar shell model},pdfauthor={Ricardo Maldonado},pdfsubject={Research draft for specialist review; not submitted}}
\usepackage[numbers,sort&compress]{natbib}
\usepackage{microtype}
\newtheorem{theorem}{Theorem}[section]
\newtheorem{proposition}[theorem]{Proposition}
\newtheorem{lemma}[theorem]{Lemma}
\newtheorem{corollary}[theorem]{Corollary}
\theoremstyle{remark}
\newtheorem{remark}[theorem]{Remark}
\newcommand{\tr}{\operatorname{tr}}
\newcommand{\Sym}{\operatorname{Sym}}
\newcommand{\ran}{\operatorname{ran}}
\newcommand{\so}{\mathfrak{so}}
\newcommand{\F}{\mathrm{F}}
\newcommand{\diag}{\operatorname{diag}}
\newcommand{\supp}{\operatorname{supp}}
\newcommand{\Lie}{\mathcal L}
\newcommand{\dd}{\mathrm d}
\newcommand{\R}{\mathbb R}
\title{Rotational linearization obstructions in a doubled Einstein--scalar shell model}
\author{Ricardo Maldonado\\\small Independent Researcher\\
\small ORCID: \href{https://orcid.org/0009-0009-3937-6527}{0009-0009-3937-6527}}
\date{15 October 2026\\\small Research draft for specialist review; not submitted}
\begin{document}
\maketitle
\begin{abstract}
We work out a rotational linearization obstruction for two Einstein--scalar
regions joined across a timelike shell whose action is a scalar-dependent
tension. The integrated momentum constraint and mixed Israel equation give
an exact functional containing the shell's area--boost-angle term. At a
rotation-invariant background its first variation vanishes for arbitrary
matching variations, including variations of the shell geometry. Consequently
a nonzero quadratic rotational functional cannot be cancelled by an
unrestricted second-order correction. In a five-dimensional radial background,
an explicit scalar master reconstruction and finite propagation produce
smooth compactly supported linear solutions before shell contact. We exhibit
an exact background with smooth interiors, including its fixed potential and shell tension,
on which a rotating quadrupole has a strictly nonzero obstruction.
For a fixed-profile quadrupole family, the six rotational conditions reduce
to commutation of two trace-free symmetric matrices. We describe the regular
and singular geometry of this finite condition, including the exact criterion
for a linearized matrix direction to integrate to a commuting curve. The
Noether and linearization-instability mechanisms are established; the result
is an explicit physical-interface application with a computable witness.
The matrix criterion is not a sufficiency theorem for nonlinear spacetimes.
\end{abstract}
\noindent\textbf{Keywords:} Einstein--scalar equations; timelike shells;
linearization instability; gravitational constraints; momentum maps.

\section{Question, contribution and relation to prior work}

A solution of the linearized field equations need not be the first derivative
of a family of exact solutions. At a symmetric background, the second
derivative of the constraints can impose a quadratic condition on a linear
solution. This is the classical linearization-stability problem
\cite{FischerMarsden1973,Moncrief1975,ArmsMarsdenMoncrief1981}.
For an internal timelike shell, applying this principle requires accounting
for the cut shared by the two bulk Cauchy sections. Compact support of the
\emph{first} variation does not remove an arbitrary \emph{second} variation
at that cut.

This paper gives a direct calculation for a pure-potential scalar shell.
Its main steps are the exact cut identity \eqref{eq:exact-functional}, the
vanishing of the first derivative on all matching corrections
\eqref{eq:first-zero}, and the explicit nonzero value
\eqref{eq:explicit-obstruction}. The witness exists on a short slab before
contact with the shell, so the central result uses a standard hyperbolic
Cauchy theorem rather than a shell initial-boundary existence theorem.
The finite-dimensional calculation in Section~\ref{sec:geometry} distinguishes
passing the rotational tests from lying in a regular part of their zero set.

The general construction is not new. Iyer--Wald supplies the bulk Noether
identity \cite{IyerWald1994}. Harlow--Wu's completed boundary Hamiltonian
already has the combination of Noether charge, boundary Lagrangian and
corner potential evaluated in Appendix~\ref{app:covariant}
\cite{HarlowWu2020}. The physical-shell Hamiltonian of H\'aj\'i\v cek--Kijowski,
with its erratum, supplies direct precedent for induced matching and the
area--boost-angle pair \cite{HajicekKijowski1998,HajicekKijowski2000Erratum}.
Recent general treatments explicitly relate second-order fields, boundary
terms and quadratic symmetry constraints
\cite{DeVuystEtAl2025,Khavkine2015}; the compact-Cauchy Taub mechanism is also
displayed directly in \cite{AltasTekin2019}.
Our calculation specializes these mechanisms to a two-sided scalar/tension
interface and evaluates a smooth obstructed family. It does not provide a
new general Noether formalism or a classification of the nonlinear constraint
range. The accompanying \href{https://github.com/maldonado-research/HDblast/blob/main/research/HDBLAST_ROTATIONAL_OBSTRUCTION_PAPER_20261015/literature/PRIMARY_COMPARISON.md}{primary-source comparison}
records the exact passages used and the remaining historical-source access gaps.

\section{Model, matching and regularity}
\label{sec:model}

The interface identity below holds in spacetime dimension $d+1$, $d\ge2$.
Use signature $(-+\cdots+)$ and $8\pi G=1$. Two bulk regions $M_+$ and $M_-$
meet at a common timelike shell $W$. Their outward spacelike unit normals
are $n_i$, and
$K_{i,ab}=h_a{}^A h_b{}^B\nabla_A n_{i,B}$.
The action is
\begin{equation}
 S=\sum_{i=\pm}\int_{M_i}\epsilon_i
 \left(\frac{R_i}{2}-\frac{(\nabla\phi_i)^2}{2}-U(\phi_i)\right)
 +\int_W\epsilon_W\,[K_++K_--\sigma(\varphi)].
 \label{eq:action}
\end{equation}
There are two outward Gibbons--Hawking--York terms and one tension term.
The smooth functions $U$ and $\sigma$ are fixed. Exact induced matching is
part of the configuration space:
\begin{equation}
 X_+^*g_+=X_-^*g_-=h,\qquad
 X_+^*\phi_+=X_-^*\phi_-=\varphi.
 \label{eq:matching}
\end{equation}
Writing $p_i=n_i\phi_i$, the metric and scalar junction residuals are
\begin{equation}
 I_{ab}=\sum_i(K_{i,ab}-K_i h_{ab})+\sigma h_{ab},\qquad
 P=p_++p_-+\sigma'(\varphi).
 \label{eq:junctions}
\end{equation}
The equations set both residuals to zero. In particular, scalar-dependent
tension has no intrinsic derivative kinetic term and its stress is
proportional to $h_{ab}$.

Let each spatial section $C_i$ be compact with common closed boundary $S$,
and assume no external spatial boundary. The induced metric $c$ on $S$ is
the same from both sides. A common reference vector $J$ is tangent to $S$
and is extended smoothly to each $C_i$. It is fixed when fields vary.
For the explicit example, $C_i$ are four-balls and $J$ is a rotation.

All differentiations below take place in a local chart of exactly matched
configurations. Moving embeddings and Cauchy sections are pulled back to
fixed reference regions by smooth collars. A hypothetical exact family is
assumed $C^2$ in its parameter in a piecewise $C^k$ norm, $k\ge4$, on a
compact slab, with uniformly nondegenerate metrics and spacelike sections.
These are sufficient classical assumptions for the stated products, traces
and differentiations. No manifold theorem for the full solution space is
assumed. The second variation may occupy every angular, metric, scalar and
shell sector compatible with exact matching.

\section{The rotational constraint with its shell cut}
\label{sec:constraint}

Let $u_i$ be the future unit normal to $C_i$, $s_i$ its outward unit normal
to $S$, and $\tau$ the common future normal to $S$ within $W$. Define the
boost angles by
\begin{equation}
 n_i=\cosh\eta_i\,s_i+\sinh\eta_i\,u_i,
 \qquad
 \tau=\cosh\eta_i\,u_i+\sinh\eta_i\,s_i.
 \label{eq:boost}
\end{equation}
For the spatial metric $\gamma_i$ and extrinsic curvature $k_i$, use the
canonical densities
\begin{equation}
 \pi_i^{ab}=\frac{\sqrt{\gamma_i}}{2}
 (k_i^{ab}-k_i\gamma_i^{ab}),\qquad
 \varpi_i=\sqrt{\gamma_i}\,u_i\phi_i,
 \qquad
 \mathcal H_{i,a}=-2D_b\pi_i^b{}_a+\varpi_iD_a\phi_i.
 \label{eq:canonical}
\end{equation}
Here $k_{ab}=\gamma_a{}^A\gamma_b{}^B\nabla_Au_B$; densities, rather than
tensors with an additional volume element, are used in the bulk integrals.

\begin{lemma}[Exact rotational cut identity]
\label{lem:cut}
For the matched configurations above,
\begin{align}
 \mathcal C_J
 &:={\sum_i}\int_{C_i}\mathcal H_{i,a}J^a
       -\int_S\sqrt c\,I(\tau,J)\nonumber\\
 &=\sum_i\int_{C_i}
       (\pi_i^{ab}\Lie_J\gamma_{i,ab}+\varpi_iJ\phi_i)
       +\int_S\sqrt c\,J(\eta_++\eta_-).
 \label{eq:exact-functional}
\end{align}
Every exact constraint-and-junction datum has $\mathcal C_J=0$.
\end{lemma}
\begin{proof}
Differentiate the orthonormal frame \eqref{eq:boost} along the cut-tangent
field $J$. Normalization and orthogonality give
\begin{equation}
 k_i(s_i,J)=-K_i(\tau,J)-J\eta_i.
 \label{eq:tilt-identity}
\end{equation}
Integration by parts of the momentum constraint gives the bulk canonical
term and the boundary term $-\int_S\sqrt c\,k_i(s_i,J)$. Since
$h(\tau,J)=0$, the trace and tension terms in \eqref{eq:junctions} have
zero mixed projection; hence $I(\tau,J)=\sum_iK_i(\tau,J)$. Substitution
of \eqref{eq:tilt-identity} proves \eqref{eq:exact-functional}.
\end{proof}

Write $\vartheta=\eta_++\eta_-$ and $\rho_S=\sqrt c$. Suppose the background
canonical fields and cut data are all $J$-invariant:
\begin{equation}
 \Lie_J\gamma_0=\Lie_J\pi_0=J\phi_0=\Lie_J\varpi_0
 =\Lie_J\rho_{S0}=J\vartheta_0=0.
 \label{eq:invariant-background}
\end{equation}
Lie derivatives of $\pi_0,\varpi_0,\rho_{S0}$ are density Lie derivatives.
These assumptions hold for the radial time-symmetric section used below.

\begin{proposition}[No linear second-order cancellation]
\label{prop:derivatives}
At a background satisfying \eqref{eq:invariant-background},
\begin{equation}
 D\mathcal C_J|_0[w]=0
 \quad\hbox{for every matching variation }w.
 \label{eq:first-zero}
\end{equation}
For a first variation $v$, the Hessian is
\begin{align}
 \mathcal T_J(v):=D^2\mathcal C_J|_0[v,v]
 &=2\sum_i\int_{C_i}
    (\delta\pi_i^{ab}\Lie_J\delta\gamma_{i,ab}
                +\delta\varpi_iJ\delta\phi_i)\nonumber\\
 &\quad+2\int_S\delta\rho_S\,J\delta\vartheta.
 \label{eq:hessian}
\end{align}
If $v$ vanishes with the required first jets near the initial shell cut,
the last integral is zero.
\end{proposition}
\begin{proof}
Differentiate \eqref{eq:exact-functional}. Terms in which the variation
multiplies a Lie derivative of a background field vanish by
\eqref{eq:invariant-background}. The remaining first-order terms are
$\int\pi_0\Lie_J\delta\gamma$,
$\int\varpi_0J\delta\phi$ and $\int\rho_{S0}J\delta\vartheta$.
Integrate them by parts. There is no exterior flux because $J$ is tangent
to the bulk cuts, and $S$ is closed. Background density invariance sets
each integral to zero. Differentiating a second time leaves exactly the
cross terms \eqref{eq:hessian}; all terms linear in a second variation
vanish by the same calculation. Because the first derivative is zero,
the displayed Hessian is independent of the choice of matching chart.
\end{proof}

\begin{theorem}[Necessary rotational continuation condition]
\label{thm:obstruction}
Let $v$ be the derivative of a classical parameter-$C^2$ family of exact
solutions of \eqref{eq:action}--\eqref{eq:junctions} through a background
satisfying \eqref{eq:invariant-background}. Then $\mathcal T_J(v)=0$.
In particular, a smooth linear solution for which this quantity is nonzero
cannot be the first derivative of such a family, even when its second
derivative is unrestricted in the bulk and on the shell.
\end{theorem}
\begin{proof}
Write the exact family in a matching chart as
$\Psi_\varepsilon=\Psi_0+\varepsilon v+
\tfrac12\varepsilon^2w+o(\varepsilon^2)$.
Since $\mathcal C_J(\Psi_\varepsilon)=0$, its second derivative is
$D^2\mathcal C_J[v,v]+D\mathcal C_J[w]=0$.
Proposition~\ref{prop:derivatives} removes the second term for every $w$.
\end{proof}

The support hypothesis removes a quadratic cut contribution of $v$.
It is not the reason the arbitrary second-order shell correction drops
out: that follows separately from \eqref{eq:first-zero}. Nor does
$\mathcal T_J=0$ prove that a correction $w$ exists. The Hessian convention
here is twice the quadratic Taylor coefficient.

\section{An admissible five-dimensional witness before shell contact}
\label{sec:witness}

Assume an actual smooth radial background of the equations exists on the
complete doubled slab, with
\begin{equation}
 g_0=\Omega(s)^2\eta,\qquad \phi_0=\Phi(s),\qquad
 s=R^2-T^2,\qquad W:\ s=b^2.
 \label{eq:background}
\end{equation}
The two copies agree, $\Omega>0$, and $f=\Phi_s\ne0$ on a neighborhood
of the annulus used below. This is an explicit hypothesis about an actual
background, not a claim that arbitrary functions $\Omega,\Phi$ solve the
nonlinear equations. A completely specified example is given in
Section~\ref{sec:explicit-background}. The registered model motivating
this calculation is documented separately \cite{MaldonadoShell2026};
its detuning and branch certification are not inputs here.

Set $\mathcal E=X^A\partial_A$ and
$L=\Omega_s/\Omega+2f_s/f$. In the region where $f\ne0$, define
\begin{equation}
 \mathscr L F=\Box_\eta F-2L\mathcal E F
       -\left(8L+\frac83sf^2\right)F.
 \label{eq:master}
\end{equation}
This is normally hyperbolic with smooth lower-order coefficients. The
reconstruction is
\begin{equation}
 h_{AB}=2F(s\eta_{AB}-3X_AX_B),\qquad
 \chi=-\frac{3(\mathcal E+4)F}{2\Omega^2f}.
 \label{eq:reconstruction}
\end{equation}

\begin{lemma}[Short-slab admissibility]
\label{lem:admissibility}
Let the smooth data $q=F|_{T=0}$ and $p=F_T|_{T=0}$ have support in
$r_1\le R\le r_2$, where $0<r_1<r_2<b$. For a sufficiently short
two-sided time slab, \eqref{eq:master} and \eqref{eq:reconstruction}
produce a smooth solution of the complete linearized bulk and shell
equations. The perturbation vanishes near the shell and near the cone
$s=0$ throughout this slab.
\end{lemma}
\begin{proof}
On $s>0$, direct substitution using the background equations gives the
operator identities
\begin{align}
 \delta\mathscr E_{AB}
 &=\Omega^{-2}(3X_AX_B-s\eta_{AB})\mathscr L F,\nonumber\\
 \delta\mathscr S
 &=-\frac{3}{2\Omega^4f}(\mathcal E+6)\mathscr L F,
 \label{eq:residual-factorization}
\end{align}
where $\mathscr E_{AB}=R_{AB}-\phi_A\phi_B-\tfrac23U g_{AB}$ and
$\mathscr S=\Box_g\phi-U'$; Appendix~\ref{app:reconstruction} supplies
the component calculation.
Extend the coefficients of the auxiliary scalar equation smoothly away
from a compact neighborhood of the intended causal support, and apply
the smooth Cauchy theorem and finite propagation
\cite[Theorem 3.2.11]{BarGinouxPfaffle2007}.
The support is contained in
$r_1-|T|\le R\le r_2+|T|$.
Choose, in addition to the background-domain restriction,
\begin{equation}
 0<T_*<\min\{r_1/2,(b-r_2)/2\}.
 \label{eq:precontact-time}
\end{equation}
Then the causal support lies in $R>|T|$ and $R<b$, separated from both
the cone and the shell $R=\sqrt{b^2+T^2}$. Shrink $T_*$ further, if needed,
so this support remains in the neighborhood where the actual background
coefficients are smooth and $f\ne0$. On this support the coefficients
are those of the actual background and
\eqref{eq:residual-factorization} applies. The reconstructed fields extend
smoothly by zero to the remaining physical region. They vanish on an open
shell neighborhood, so zero embedding variation satisfies every linearized
matching and junction condition there. All initial bulk constraints follow
from the bulk equations. The auxiliary coefficient extension has made no
assertion about a nonlinear background outside the physical neighborhood.
\end{proof}

The first-order embedding and field jets vanish near the shell. Thus the
pullback identifications for a candidate exact family can be chosen with
zero first parameter jet near the entire initial Cauchy section, including
the time jet needed for its canonical data. This preserves the explicit
first variation below while allowing arbitrary second-order pullback terms.
At $T=0$, write $D=r\partial_r$. Direct ADM variation gives
\begin{equation}
 \begin{aligned}
 \delta\gamma_{ij}&=2q(r^2\delta_{ij}-3x_ix_j),&
 \delta\pi^{ij}&=-\frac{3p}{2\Omega}x_ix_j,\\
 \delta\phi&=-\frac{3(D+4)q}{2\Omega^2f},&
 \delta\varpi&=-\frac{3\Omega(D+5)p}{2f}.
 \end{aligned}
 \label{eq:initial-reconstruction}
\end{equation}
The shift by five in the last formula follows from
$[\partial_T,\mathcal E]=\partial_T$.
Define
\begin{equation}
 \mathcal B(q,p)=\int\dd r\,\dd\Omega_3
 \left[m_r(D+4)q(D+5)p+n_rqp\right],\quad
 m_r=\frac{9r^3}{4\Omega f^2},\quad n_r=\frac{6r^7}{\Omega}.
 \label{eq:bilinear}
\end{equation}
Use identical data on the two copies (or opposite signs on both fields
of one copy). Substitution into \eqref{eq:hessian} yields
\begin{equation}
 \mathcal T_J(v)=4\mathcal B(Jq,p).
 \label{eq:master-charge}
\end{equation}

Choose $Q\in C_c^\infty(I)$ nonzero, $I\Subset(0,b)$ in the regular
annulus, a real harmonic $H$ with $JH\ne0$, and a nonzero real inverse-time
constant $\omega$. Put $q=QH$, $p=\omega QJH$. Then
\begin{equation}
 \mathcal T_J(v)=4\omega\mathcal B_{\rm rad}(Q,Q)
                              \int_{S^3}(JH)^2\dd\Omega_3.
 \label{eq:charge-separation}
\end{equation}
No symmetry of the radial bilinear is assumed. Compact radial integration
by parts gives
\begin{equation}
 \mathcal B_{\rm rad}(Q,Q)
 =\int_I\left[m_r(DQ)^2+
       \{m_r(2+9r^2L)+n_r\}Q^2\right]\dd r.
 \label{eq:radial-diagonal}
\end{equation}
If $w_I$ is the length of $I$, let
$m_* =\min_{\bar I}(m_rr^2)>0$ and
$M_* =\max_{\bar I}m_r(-2-9r^2L)_+$. The one-dimensional Poincar\'e
inequality implies
\begin{equation}
 \mathcal B_{\rm rad}(Q,Q)\ge
 \left(\frac{\pi^2m_*}{w_I^2}-M_*\right)\int_IQ^2\dd r>0
 \label{eq:radial-positive}
\end{equation}
on a sufficiently thin annulus. This is positivity of one diagonal radial
form, not positivity of a gravitational energy.

For $n=x/r$, take $H=n_1^2-n_2^2$ and
$J=x_1\partial_2-x_2\partial_1$. Then
$JH=-4n_1n_2$ and $\int_{S^3}(JH)^2=4\pi^2/3$, so
\begin{equation}
 \boxed{\mathcal T_J(v)=\frac{16\pi^2}{3}
                  \omega\mathcal B_{\rm rad}(Q,Q)\ne0.}
 \label{eq:explicit-obstruction}
\end{equation}
Lemma~\ref{lem:admissibility} makes this an actual smooth linear solution.
Theorem~\ref{thm:obstruction} excludes a classical parameter-$C^2$ nonlinear
continuation with this tangent. Here an admissible pure-gauge variation is
the infinitesimal joint diffeomorphism action on bulk fields and shell
embeddings, generated by a smooth vector field with a local flow near the
compact initial section. Such a variation is tangent to an exact
diffeomorphic family; the nonzero obstruction therefore excludes it.

\subsection{A fully specified background and witness}
\label{sec:explicit-background}

In a fixed length unit, set $z=1+s$ and choose
\begin{equation}
 \Omega=z^{-2},\qquad \Phi=\sqrt6\log z,\qquad
 U(\phi)=84e^{2\phi/\sqrt6}-36e^{3\phi/\sqrt6}.
 \label{eq:explicit-background}
\end{equation}
These fields solve the bulk equations on $z>0$. Indeed, writing
$u=\log\Omega$, the conformal Ricci formula gives
\begin{align}
 R_{AB}
 &=12(u_s^2-u_{ss})X_AX_B
   +(-16u_s-4su_{ss}-12su_s^2)\eta_{AB}\nonumber\\
 &=\frac{24X_AX_B+(32-24s)\eta_{AB}}{z^2}
  =\Phi_A\Phi_B+\frac23 U(\Phi)g_{AB},
 \label{eq:explicit-ricci}
\end{align}
because $\Phi_s=\sqrt6/z$ and $U(\Phi)=(48-36s)z^2$.
Independently,
\begin{equation}
 \Box_g\Phi=\Omega^{-2}(4s\Phi_{ss}+10\Phi_s+12su_s\Phi_s)
 =\sqrt6(10-18s)z^2=U'(\Phi).
 \label{eq:explicit-scalar}
\end{equation}

Join two identical interiors $s\le b^2$ at $b=1/2$.
Their outward normal is $n^A=X^A/(b\Omega_b)$ and their extrinsic
curvatures and scalar normal derivatives are
\begin{equation}
 K_{i,ab}=a_bh_{ab},\quad
 a_b=\frac{(1-3b^2)(1+b^2)}{b}=\frac58,\qquad
 p_b=2\sqrt6\,b(1+b^2)=\frac{5\sqrt6}{4}.
\end{equation}
With $\phi_b=\sqrt6\log(5/4)$, the fixed affine tension
\begin{equation}
 \sigma(\phi)=\frac{15}{4}-\frac{5\sqrt6}{2}(\phi-\phi_b)
 \label{eq:explicit-tension}
\end{equation}
satisfies $\sigma(\phi_b)=6a_b$ and $\sigma'(\phi_b)=-2p_b$.
Thus both residuals in \eqref{eq:junctions} vanish exactly. Induced
matching is automatic. On $|T|\le1/2$, the interior has
$-1/4\le s\le1/4$, so the Cartesian fields are smooth through the center
and cone within each interior, with $\Omega>0$ and $f\ne0$.
Across the material shell the metric is piecewise smooth, with the
extrinsic-curvature jump fixed by the junction equations. The initial section is time
symmetric and rotation invariant.

For this background,
\begin{equation}
 L=-\frac4{1+s},\qquad m_r=\frac38r^3(1+r^2)^4,\qquad
 n_r=6r^7(1+r^2)^2,
\end{equation}
and the zeroth-order radial coefficient is
\begin{equation}
 m_r(2+9r^2L)+n_r
 =m_r\frac{2-32r^2-18r^4}{(1+r^2)^2}>0
 \quad(0<r\le1/5).
 \label{eq:explicit-positive}
\end{equation}
The numerator decreases on this interval and equals $432/625$ at its
upper endpoint. Hence any nonzero $Q\in C_c^\infty((1/10,1/5))$ has
$\mathcal B_{\rm rad}(Q,Q)>0$, without the thin-annulus estimate.
For example, take $Q=\exp[-1/((r-11/100)(19/100-r))]$ between
$11/100$ and $19/100$, and zero elsewhere. Take $T_*=1/32$ and
$\omega=1$ in the fixed units. Together with the displayed $H$ and $J$,
these choices specify a smooth linear state with
\eqref{eq:explicit-obstruction} strictly positive on a fully specified
exact background.

This illustrative potential and tension differ from the registered
HDBLAST model. The example establishes that the present hypotheses are
realized; it makes no priority claim for the exact background and no
claim about its long-time stability or cosmological interpretation.

\section{Quadrupole charges and commuting matrices}
\label{sec:quadrupoles}

Keep $Q$, $\omega$ and the copy parity fixed, with
$\mathcal B_{\rm rad}(Q,Q)>0$. Every real degree-two harmonic on $S^3$
is $H_A(n)=n^TAn$ with $A\in V:=\Sym_0(4,\R)$. The family
\begin{equation}
 q=QH_A,\qquad p=\omega QH_B,\qquad (A,B)\in V\times V
 \label{eq:quadrupole-family}
\end{equation}
has dimension eighteen. For $M^T=-M$ and $J_M=(Mn)\cdot\partial_n$,
the exact sphere moments give
\begin{equation}
 \int_{S^3}H_AH_B=\frac{\operatorname{vol}(S^3)}{12}\tr(AB),
 \qquad J_MH_A=H_{[A,M]}.
 \label{eq:sphere-moments}
\end{equation}
Consequently
\begin{equation}
 \mathcal T_{J_M}(v)=
 \frac{\operatorname{vol}(S^3)}{3}\omega
 \mathcal B_{\rm rad}(Q,Q)\tr\bigl(B[A,M]\bigr).
 \label{eq:matrix-charge}
\end{equation}
The six rotational tests therefore vanish exactly when $[A,B]=0$.
For $J_{ij}=x_i\partial_j-x_j\partial_i$, the coefficient is
$-(2\operatorname{vol}(S^3)/3)\omega\mathcal B_{\rm rad}(Q,Q)[A,B]_{ij}$.
This calculation concerns the restricted family
\eqref{eq:quadrupole-family}, not every perturbation or every adjoint
constraint mode.

\section{Regular and singular quadrupole test states}
\label{sec:geometry}
For the finite quadrupole condition, consider the general map
\[
 V=\Sym_0(n,\mathbb R),\qquad
 \mu:V\times V\longrightarrow\so(n),\qquad \mu(A,B)=[A,B],
 \qquad Z=\mu^{-1}(0),\qquad n\geq2.
\]
The metrics are the Frobenius metric on matrices and its product metric on
$V\times V$.  This is standard commuting symmetric-matrix geometry.  Under
the canonical identification $T^*V\simeq V\times V$, the commutator represents
the moment map for simultaneous orthogonal conjugation, up to the sign fixed
by the dual-pairing convention.  The claims below do not concern the range
of a gravitational constraint operator, interface conditions, or nonlinear
spacetime continuation.  The momentum-map context is classical \cite{ArmsMarsdenMoncrief1981}.

Fix $(A,B)\in Z$.  Simultaneous orthogonal diagonalization gives the
\emph{maximal} joint eigenspace decomposition
\[
 \mathbb R^n=\bigoplus_{s=1}^r E_s,\qquad
 A|_{E_s}=\alpha_s I,\qquad B|_{E_s}=\beta_s I,
 \qquad (\alpha_s,\beta_s)\ne(\alpha_t,\beta_t)\quad(s\ne t).
\]
Write $m_s=\dim E_s$, let $P_s$ be its orthogonal projector, and put
$a_{st}=P_s a|_{E_t}$ and $a_{\mathrm d}=\sum_sP_saP_s$, with analogous
notation for $b$.  In a simultaneous orthonormal eigenbasis, use repeated
lists $\alpha_i,\beta_i$ and define
\[
 c_{ij}=\beta_j-\beta_i,\qquad d_{ij}=\alpha_i-\alpha_j,
 \qquad g_{ij}^2=c_{ij}^2+d_{ij}^2\quad(i<j).
\]

\begin{proposition}[Rank, optimal inverse, and regularity]\label{prop:inverse}
The derivative $L=D\mu_{(A,B)}$ is
\[
 L(a,b)=[a,B]+[A,b],\qquad L(a,b)_{ij}=c_{ij}a_{ij}+d_{ij}b_{ij}.
\]
Its nonzero singular values are exactly the positive $g_{ij}$.  Equivalently,
for each $s<t$ the singular value
\[
 \sqrt{(\alpha_s-\alpha_t)^2+(\beta_s-\beta_t)^2}
\]
occurs with multiplicity $m_s m_t$.  Hence, with
$d=\dim V=n(n+1)/2-1$ and $N=\dim\so(n)=n(n-1)/2$,
\[
 \operatorname{rank} L=N-\sum_s\binom{m_s}{2},\qquad
 \dim\ker L=2d-N+\sum_s\binom{m_s}{2}.
\]
The range comprises precisely the skew matrices $S$ with $P_sSP_s=0$ for
all $s$.  The Moore--Penrose inverse $R=L^\dagger$ is the pair of symmetric
zero-diagonal matrices with
\[
 (RS)^{(1)}_{ij}=\frac{c_{ij}S_{ij}}{g_{ij}^2},\qquad
 (RS)^{(2)}_{ij}=\frac{d_{ij}S_{ij}}{g_{ij}^2}\qquad(g_{ij}>0),
\]
and both entries zero when $g_{ij}=0$.  Here the two displayed components
are extended symmetrically.  Then $LR$ is the orthogonal range projection.
If at least one $g_{ij}>0$,
\[
 \|R\|_{\so(n)\to V\times V}
   =\frac{1}{\min_{g_{ij}>0}g_{ij}}.
\]
For $L=0$, $R=0$ and $\|R\|=0$.  In particular, a right inverse on all of
$\so(n)$ exists exactly when every joint eigenvalue is simple.  In that
case $R$ is the unique solution of minimum product norm for each target,
and its operator norm is the smallest possible among all linear right
inverses.
\end{proposition}

\begin{proof}
For $i<j$, the normalized symmetric and skew matrix units are respectively
$(E_{ij}+E_{ji})/\sqrt2$ and $(E_{ij}-E_{ji})/\sqrt2$.  On the two normalized
input directions belonging to this edge, $L$ is the single row
$(c_{ij},d_{ij})$.  Distinct edges are orthogonal; all diagonal input
directions are in the kernel.  The trace restrictions affect only those
diagonal inputs.  Thus each edge contributes one singular value $g_{ij}$
unless it lies within a joint eigenspace.  This proves the singular-value,
rank, and range statements.  The inverse formula is the minimum-norm
solution of the scalar equation $c_{ij}a_{ij}+d_{ij}b_{ij}=S_{ij}$.
Consequently,
\[
 \|RS\|_{V\times V}^2
   =2\sum_{i<j:\,g_{ij}>0}\frac{S_{ij}^2}{g_{ij}^2}.
\]
The norm formula follows by concentrating $S$ on an edge of minimum
positive gap.  Every solution for that target has norm at least the
minimum-norm solution, which also proves optimality among right inverses.
There is no extra factor $\sqrt2$ in the singular values: the input and
output Frobenius normalizations agree.
\end{proof}

\begin{theorem}[All velocities of commuting curves]\label{thm:velocities}
For $a,b\in V$, the following conditions are equivalent.
\begin{enumerate}
\item There is a $C^1$ curve $t\mapsto(A(t),B(t))\in Z$ through $(A,B)$
with derivative $(a,b)$ at $0$.
\item $L(a,b)=0$ and $[a_{ss},b_{ss}]=0$ on every maximal joint eigenspace
$E_s$.
\item $L(a,b)=0$ and there are unrestricted $c,d\in V$ satisfying the
second-order matrix equation
\[
 L(c,d)=-2[a,b].
\]
\item There is a real-analytic such curve with derivative $(a,b)$ at $0$.
\end{enumerate}
For a velocity satisfying these conditions, define a skew matrix $K$ by
$K_{ss}=0$ and, for $s\ne t$,
\[
 K_{st}=
 \frac{(\alpha_t-\alpha_s)a_{st}+(\beta_t-\beta_s)b_{st}}
 { (\alpha_t-\alpha_s)^2+(\beta_t-\beta_s)^2 }.
\]
One explicit realization is
\begin{align*}
 Q(t)&=(I+tK/2)(I-tK/2)^{-1},\\
 A(t)&=Q(t)(A+t a_{\mathrm d})Q(t)^{\mathsf T},&
 B(t)&=Q(t)(B+t b_{\mathrm d})Q(t)^{\mathsf T}.
\end{align*}
This rational curve is defined and real analytic for every real $t$.
\end{theorem}

\begin{proof}
First, $L(a,b)=0$ is necessary by differentiation.  Between distinct joint
blocks it reads
\[
 (\beta_t-\beta_s)a_{st}-(\alpha_t-\alpha_s)b_{st}=0.
\]
The coefficient pair is nonzero.  Therefore the displayed definition of
$K$ gives
\[
 a_{st}=(\alpha_t-\alpha_s)K_{st},\qquad
 b_{st}=(\beta_t-\beta_s)K_{st}\quad(s\ne t).
\]
Symmetry of $a,b$ and reversal of the gaps imply $K_{ts}=-K_{st}^{\mathsf T}$.
In particular,
\begin{equation}\label{eq:offblock-cancellation}
 a_{st}b_{ts}-b_{st}a_{ts}=0\quad(s\ne t),\qquad
 P_s[a,b]P_s=[a_{ss},b_{ss}].
\end{equation}

For necessity of the block condition, differentiability at $0$ alone is
enough.  Put $X(t)=A(t)-A=ta+o(t)$ and $Y(t)=B(t)-B=tb+o(t)$.  Every
commutator with $A$ or $B$ has zero diagonal joint block, so the exact
equation $[A+X(t),B+Y(t)]=0$ gives
\[
 0=P_s[X(t),Y(t)]P_s=t^2P_s[a,b]P_s+o(t^2).
\]
Divide by $t^2$ and use \eqref{eq:offblock-cancellation}.  This proves
$(1)\Rightarrow(2)$ without assuming a second derivative of the curve.

For sufficiency, the block condition says $[a_{\mathrm d},b_{\mathrm d}]=0$.
Both block-diagonal matrices commute with $A$ and $B$, since the latter
are scalar on each joint block.  Thus
$[A+t a_{\mathrm d},B+t b_{\mathrm d}]=0$ identically.  For real $t$,
$I-tK/2$ is invertible: if $(I-tK/2)x=0$, taking the real inner product with
$x$ gives $\|x\|^2=0$.  The Cayley matrix $Q(t)$ is orthogonal, has
$Q(0)=I$, and has $Q'(0)=K$.  Hence the proposed curve is symmetric,
tracefree, real analytic, and commuting, and
\[
 A'(0)=a_{\mathrm d}+[K,A]=a,\qquad
 B'(0)=b_{\mathrm d}+[K,B]=b.
\]
This proves $(2)\Rightarrow(4)\Rightarrow(1)$.

Finally, Proposition~\ref{prop:inverse} says the second-order equation is
solvable exactly when $P_s[a,b]P_s=0$ for all $s$.  Under $L(a,b)=0$,
\eqref{eq:offblock-cancellation} identifies this with the block condition.
This proves $(2)\Leftrightarrow(3)$.  Explicit accelerations of the curve are
\[
 c=2[K,a_{\mathrm d}]+[K,[K,A]],\qquad
 d=2[K,b_{\mathrm d}]+[K,[K,B]],
\]
and $R(-2[a,b])$ is the correction of minimum product norm.
\end{proof}

\begin{remark}
The condition is \emph{not} $[a,b]=0$ on all of $\mathbb R^n$.  Off-block
components of this commutator can be canceled by second-order corrections.
For example, let $A=\diag(1,-1,0,0)$, $B=0$,
$a=E_{12}+E_{21}$, and $b=A$.  Then $L(a,b)=0$ and every diagonal joint
block commutator is zero, although $[a,b]\ne0$.  Splitting a repeated
maximal joint block into artificial one-dimensional blocks would instead
discard genuine obstructions and is invalid.
\end{remark}

\begin{corollary}[Smooth points and a simple pencil]\label{cor:regular}
The following are equivalent at a commuting pair: $L$ is surjective; all
maximal joint eigenspaces are one dimensional; some real pencil $A+\tau B$
has simple spectrum; and $Z$ is a $C^1$ embedded submanifold in a
neighborhood of that pair.  Its dimension there is
\[
 2\dim V-\dim\so(n)=\frac{n^2+3n-4}{2}.
\]
Neither $A$ nor $B$ need individually have simple spectrum.
\end{corollary}

\begin{proof}
Surjectivity is Proposition~\ref{prop:inverse}, and the implicit-function
theorem gives a smooth submanifold at every such point.  If $m_s\geq2$,
choose an internal tracefree diagonal matrix $a$ with entries $1,-1$ and
$b=E_{ij}+E_{ji}$ on those two coordinates.  The velocities $(a,0)$ and
$(0,b)$ are realized by straight commuting curves, while their sum fails
Theorem~\ref{thm:velocities}.  The velocity set of a $C^1$ embedded
submanifold would be a vector space.  Thus no repeated joint block is a
smooth point.  For distinct joint pairs, a collision
$\alpha_i+\tau\beta_i=\alpha_j+\tau\beta_j$ excludes at most one value of
$\tau$ per pair; if the two $\beta$ values agree, no value is excluded.
Avoiding this finite set gives a simple pencil.  A repeated joint block
prevents a simple pencil for every $\tau$.
\end{proof}

\noindent
A local polynomial chart at any regular pair can therefore be based on a
simple pencil rather than on $A$ alone.  Fix a suitable $\tau$ and set
$C=A+\tau B$, $D=B$.  Every symmetric tracefree $D$ commuting with a
nearby simple $C$ is uniquely
\[
 D=\sum_{k=1}^{n-1}c_k\left(C^k-\frac{\tr(C^k)}{n}I\right),\qquad
 A=C-\tau D,\quad B=D.
\]
Interpolation on the $n$ distinct eigenvalues proves existence and
uniqueness.  The coefficients depend smoothly on $(C,D)$, for example by
inverting the nonsingular Frobenius Gram matrix of the displayed
polynomial basis.  The $\dim V+n-1$ parameters agree with
Corollary~\ref{cor:regular}.  In dimension four, an example with neither
individual matrix simple is
\[
 A=\diag(-1,-1,1,1),\qquad B=\diag(-1,1,-1,1).
\]
All four joint pairs are distinct, and $A+2B=\diag(-3,1,-1,3)$ is simple.

\paragraph{Dimension-four strata and the audited seeds.}
For $n=4$, the complete multiplicity list gives
\[
\begin{array}{c|r|r}
\text{joint multiplicities}&\operatorname{rank}L&\dim\ker L\\ \hline
1,1,1,1&6&12\\
2,1,1&5&13\\
2,2&4&14\\
3,1&3&15\\
4&0&18
\end{array}
\]
Every row other than the first has a nonlinear set of realizable velocities.
At the norm-two seed $A=B=\diag(1,-1,0,0)$, take
$a=\diag(0,0,1,-1)$ and $b=E_{34}+E_{43}$.  Their sum satisfies the linear
equation but has $[a,b]_{34}=2$ in the repeated block, so no $C^1$ commuting
curve realizes that velocity.  This proves the seed's genuine
singularity, rather than merely a failure of the chosen equations to have
independent derivatives.

For the norm-matched replacement
\[
 A=B=\frac1{\sqrt{10}}\diag(3,1,-1,-3),\qquad
 \tr(A^2)=\tr(B^2)=2,
\]
the minimal single-matrix spectral gap is $\gamma=2/\sqrt{10}$ and the
minimal joint gap is $\sqrt2\gamma=2/\sqrt5$.  Thus the inverse bound is sharp:
\[
 \|R\|=\frac{\sqrt5}{2}.
\]
When $A=B$ has distinct eigenvalues $\lambda_i$, the general inverse
reduces to $a=-C/2$, $b=C/2$, where
$C_{ij}=S_{ij}/(\lambda_i-\lambda_j)$ and $C_{ii}=0$.

\paragraph{Scope for the quadrupole application.}
When the restricted charge equations have independently been established
to reduce to $[A,B]=0$ with a nonzero prefactor, the preceding statements
classify their entire finite curve-velocity obstruction and the optimal
linear inverse.  They make no claim that charge-zero data, or finite
second-order solvability, suffice for the Einstein--scalar bulk and shell
equations.  The linear map $D\mu$ and its range must not be substituted for
the full constraint/interface operator.


\section{Scope, reproducibility and next questions}

The obstruction is a necessary condition for a specified classical exact
family. It neither constructs a nonlinear family through a charge-zero
state nor proves second-order solvability, stability or a positive physical
energy. The matrix curve result is sufficient only for the finite equation
$[A,B]=0$. It cannot be substituted for solving the full Einstein--scalar
constraint and interface operator. A simple matrix spectrum removes one
finite-dimensional degeneracy; it does not supply a PDE right inverse.

The direct witness here remains away from the shell during the short slab.
It therefore proves the obstruction without a contact theorem. The earlier
technical checkpoint contains a separate finite-angular contact construction
\cite{MaldonadoShell2026}; its stronger domain and compatibility hypotheses
are not implicit in Lemma~\ref{lem:admissibility}. The present paper does not
infer a nonzero jet at the exact causal front.

A constructive nonlinear continuation would require identifying the range
of the full linearized constraint and interface operator, the gauge and trace domains,
and a second-order correction with all generated angular sectors. Further
orders and estimates of the nonlinear remainder would then remain. The
calculation makes no assertion about an observational scale or the origin
of the Big Bang.

The accompanying source package contains the derivation, explicit finite
matrix checks, deliberately incorrect-formula controls, input identities,
and the primary-source comparison. The checks test algebra and selected
exact examples. They are not a computer-certified analytic PDE proof or
external peer review. Published earlier checkpoints are preserved as
separate records; the present manuscript is a focused review draft.
The \href{https://github.com/maldonado-research/HDblast/tree/main/research/HDBLAST_ROTATIONAL_OBSTRUCTION_PAPER_20261015}{draft and reproducible supplement}
are available in the public research repository.

\paragraph{AI assistance and author review.}
Substantial OpenAI Codex assistance was used in derivation, programming,
literature comparison and drafting, including internal reviews by separately
delegated GPT-6 Astra and GPT-6.1 Sol agents. These are AI-assisted checks,
not independent human peer review. The human author is responsible for
checking the mathematics, sources, interpretation and final disclosure
before any submission. Funding, conflicts, contact details and all journal
declarations must be confirmed by the author for the submitted version.

\appendix
\section{Bulk reconstruction calculation}
\label{app:reconstruction}

On $s>0$, set $r=\sqrt s$, $\dd y=\Omega\dd r$,
$\rho=r\Omega$, $p_b=\phi_y=2rf/\Omega$ and $a=\rho_y/\rho$.
Here $p_b$ denotes the radial background scalar derivative, not the master
velocity $p$ of Section~\ref{sec:witness}. The background equations are
\begin{equation}
 (p_b)_y=U_\phi-4ap_b,\quad
 a_y=-\rho^{-2}-p_b^2/3,\quad
 a^2=\rho^{-2}+p_b^2/12-U/6.
\end{equation}
For $P=r^4F$, define
\begin{equation}
 \mathcal M[P]=P_{yy}-(2a+2(p_b)_y/p_b)P_y
  +\rho^{-2}(\Box_{\mathrm{dS}_4}+4)P-\tfrac23p_b^2P.
\end{equation}
The ambient wave operator is
$\Box_\eta=\partial_r^2+4r^{-1}\partial_r+
r^{-2}\Box_{\mathrm{dS}_4}$. Differentiation gives
$\mathcal M[r^4F]=\Omega^{-2}r^4\mathscr L F$.
With $\Psi=P/\rho^2$ and $\chi=-3P_y/(\rho^2p_b)$,
the reconstructed metric is
$h_{yy}=-4\Psi$, $h_{y\mu}=0$,
$h_{\mu\nu}=2\rho^2\Psi\gamma_{\mu\nu}$.
For arbitrary spacetime functions $\Psi$, its Ricci variation is
\begin{align}
 \delta R_{yy}&=-4\Psi_{yy}-16a\Psi_y+2\rho^{-2}\Box_\gamma\Psi,\nonumber\\
 \delta R_{y\mu}&=-3D_\mu(\Psi_y+2a\Psi),\nonumber\\
 \delta R_{\mu\nu}&=\gamma_{\mu\nu}
 \left[-\Box_\gamma\Psi-\rho^2\{\Psi_{yy}+10a\Psi_y
                      +6(a_y+4a^2)\Psi\}\right].
\end{align}
Put
$C=p_b\chi_y-(p_b)_y\chi-3\rho^{-2}(\Box_\gamma+4)\Psi
+2p_b^2\Psi$.
Substitution gives
\begin{equation}
 \delta\mathscr E_{y\mu}=0,\quad
 \delta\mathscr E_{yy}=-2C/3,\quad
 \delta\mathscr E_{\mu\nu}=\rho^2\gamma_{\mu\nu}C/3,\quad
 C=-3\mathcal M[P]/\rho^2.
\end{equation}
The scalar variation satisfies
$p_b\delta\mathscr S=(\partial_y+4a)C$, giving
\eqref{eq:residual-factorization}. The pre-contact support stays in this
$s>0$ patch, so no analytic continuation of an operator identity through
$s=0$ is used in the proof.

\section{Covariant consistency of the interface cancellation}
\label{app:covariant}

This appendix records the local boundary calculation behind the same
physical interpretation; the direct rotational proof does not depend on a
general phase-space existence theorem. Use the Einstein--scalar bulk
potential, with $q_{AB}=\delta g_{AB}$, $q=g^{AB}q_{AB}$ and
$q^{AB}=g^{AC}g^{BD}q_{CD}$:
\begin{equation}
 \theta^A=\sqrt{-g}\left[\tfrac12(\nabla_Bq^{AB}-\nabla^Aq)
                         -\nabla^A\phi\,\delta\phi\right].
\end{equation}
Represent the two outward boundary forms in the same co-oriented shell
density convention, including the Noether potentials and corners.
For $v_i^a=h^{ab}n_i^A\delta g_{i,Ab}$, let
$\beta_i=-\tfrac12i_{v_i}\epsilon_W$ and
$\ell=(K_++K_--\sigma)\epsilon_W$.
The unrestricted GHY variation gives
\begin{equation}
 \sum_i\theta_i|_W+\delta\ell
 =-\epsilon_W\left(\tfrac12I^{ab}\delta h_{ab}
                         +P\delta\varphi\right)+\dd_W\beta,
 \qquad \beta=\beta_++\beta_-.
 \label{eq:boundary-variation}
\end{equation}
Normal and mixed metric variations have been retained.
For a fixed common shell-tangent vector $\xi$, let
$Q_{\xi,i}=-\tfrac12*_i\dd\xi_i^\flat$. Set
$A_i^a=h^a{}_Bn_i^A\nabla_A\xi_i^B$. Then
\begin{equation}
 Q_{\xi,i}|_W=-\tfrac12i_{A_i+K_i\cdot\xi}\epsilon_W,
 \qquad
 \beta_i(\Lie_\xi g)=-\tfrac12i_{A_i-K_i\cdot\xi}\epsilon_W.
\end{equation}
The arbitrary normal extension $A_i$ cancels, leaving
\begin{equation}
 B_\xi:=\sum_iQ_{\xi,i}|_W+i_\xi\ell-\beta(\Lie_\xi)
       =-i_{I\cdot\xi}\epsilon_W.
 \label{eq:boundary-charge}
\end{equation}
The combination on the left is the established completed boundary charge
\cite[Eq.~(2.51)]{HarlowWu2020}; its explicit Israel-residual value is the
interface calculation here. Its vanishing uses the physical junction
equation, not a fixed induced metric at an external wall.

For moving embeddings, exact pullback gives
$\delta(Y^*g)=Y^*(\delta g+\Lie_\chi g)$, where
$\chi=\delta Y\circ Y^{-1}$ \cite{Speranza2018}. Thus the second variation
contains $\ddot g+2\Lie_\chi\dot g+\Lie_{\dot\chi}g+\Lie_\chi^2g$.
Equation~\eqref{eq:matching} concerns the full pullback, including these
terms. The matching chart used in Section~\ref{sec:constraint} retains them
and the normalization of the shell normals. No vanishing of an unrestricted
second-order graph has been assumed.

\bibliographystyle{unsrtnat}
\bibliography{REFERENCES}
\end{document}
